% Generated draft: requires mathematical and source review.
\documentclass[11pt]{article}
\usepackage[margin=25mm]{geometry}
\usepackage{fontspec}
\setmainfont{Latin Modern Roman}
\usepackage{amsmath,amssymb,mathtools,unicode-math}
\setmathfont{Latin Modern Math}
\usepackage{enumitem,needspace,xurl,hyperref}
\setlength{\parindent}{0pt}
\setlength{\parskip}{5pt}
\setlength{\emergencystretch}{3em}
\newcommand{\N}{\mathbb{N}}
\newcommand{\Z}{\mathbb{Z}}
\newcommand{\Q}{\mathbb{Q}}
\newcommand{\R}{\mathbb{R}}
\newcommand{\C}{\mathbb{C}}
\newcommand{\F}{\mathbb{F}}
\newcommand{\A}{\mathbb{A}}
\newcommand{\PP}{\mathbb{P}}
\newcommand{\E}{\mathbb{E}}
\newcommand{\eps}{\varepsilon}
\newcommand{\dd}{\,\mathrm{d}}
\begin{document}
\section*{UnsolvedMath ID: 30004140}
\textbf{Source problem code:} catalogue entry
\subsection{Definitions and mathematical statement}
Let \(X=\R^n/\Lambda\) be a tropical torus as in the supplied record. Here \(\Lambda\) denotes the period lattice. The record does not specify its compatibility with an integral affine structure, the permitted dimensions, or the cohomology and positivity conventions to be used.

Once these structures and conventions have been fixed, write \(H^{1,1}(X;\Z)\) for the group of integral classes of bidegree \((1,1)\). In the usual codimension-one interpretation, an effective tropical cycle is a balanced weighted polyhedral cycle of pure codimension one with nonnegative integral weights. Provided the chosen theory supplies a cycle-class map, define
\[
\mathcal E_X=\bigl\{\alpha\in H^{1,1}(X;\Z):\alpha=[Z]\text{ for an effective tropical cycle }Z\bigr\}.
\]
This notation is conditional on those choices; it does not resolve the missing conventions in the source.

The supplied problem asks for the following three investigations.
\begin{enumerate}
\item Determine whether the stated positivity criterion for membership in \(\mathcal E_X\) admits a converse. If \(\mathcal P_X\) denotes the integral classes satisfying the intended positivity condition, the proposed characterization can be recorded schematically as
\[
\mathcal E_X\stackrel{?}{=}\mathcal P_X.
\]
The record does not define \(\mathcal P_X\), state the original implication, or specify whether positivity is strict, weak, or semidefinite. Thus this display identifies the intended converse question but is not a fully specified conjecture.
\item Determine whether tropical analogues of the Nakai--Moishezon and Kleiman criteria exist. A precise statement requires the relevant class of tropical spaces, the notion of ampleness or positivity, the intersection pairing, the test cycles, and the numerical cones and their topology. These data, including whether the proposed criteria concern only tropical tori, are absent from the record.
\item Determine which combinatorial types of tropical hypersurfaces permit large or small ranges of Picard ranks. A precise formulation requires a specified family of hypersurfaces, a definition of combinatorial type, a convention for Picard rank, and a quantified meaning of a large or small range. The record also does not specify whether these hypersurfaces lie in \(X\) or in another ambient tropical space.
\end{enumerate}
Consequently, the supplied material specifies a collection of research questions, but does not determine a unique theorem or conjecture with complete hypotheses and quantifiers.
\subsection{Short English statement}
The problem asks whether positivity characterizes effective integral \((1,1)\)-classes on tropical tori, whether tropical versions of the Nakai--Moishezon and Kleiman criteria exist, and how hypersurface combinatorics constrain Picard ranks. The positivity condition, numerical framework, hypersurface families, and rank conventions are unspecified, so a fully precise conjecture requires further source clarification.
\subsection{Sources}
\begin{itemize}
\item Supplied local UnsolvedMath registry record, permanent ID \(30004140\), title \emph{Positivity Criteria for Tropical Tori}, source problem code \texttt{catalogue entry}; local registry: \texttt{data/unsolvedmath/problems.json}. The supplied collection is \emph{Oberwolfach Reports Open Problems}. No individual proposer is identified; the supplied proposed year is 2019.
\item Supplied report citation: \emph{Tropical Geometry: new directions (2020)}, Oberwolfach Reports, pp.\ 1245--1308, DOI \texttt{10.4171/owr/2019/21}. Supplied source link: \url{https://doi.org/10.4171/owr/2019/21}. The background describes Oberwolfach Reports as documenting mathematical workshops and their open-problem discussions. The report text has not been independently consulted here.
\item Supplied literature source, labelled \emph{Tropical geometry problem report}: \url{https://ems.press/content/serial-article-files/46800}.
\item Literature assessment dated exactly \texttt{2026-08-22}: the supplied curation labels the questions open, describes positivity, tropical Nakai--Moishezon or Kleiman criteria, and Picard-rank criteria as open directions, and reports that the literature checked provides no general theorem resolving them. This is supplied dated curation derived from the OWR material, not an independently verified literature assessment. The background explicitly requests independent verification of time-sensitive status and cited claims.
\item The registry labels the statement \texttt{corrected\_verified} and reports a manual semantic, typographic, and LaTeX review claiming self-containment and preservation of the mathematical target. This is supplied review metadata; it does not resolve the formulation gaps listed here or establish mathematical correctness.
\end{itemize}
\end{document}
